% =====================================================================
%  PREÁMBULO — Análisis DAR
%  Configuración de paquetes, colores, tablas y leyendas.
% =====================================================================

% ---------- Codificación e idioma -----------------------------------
\usepackage[utf8]{inputenc}
\usepackage[T1]{fontenc}
% Si tu distribución tiene el idioma instalado (Overleaf/TeX Live full),
% descomenta la línea siguiente para partición silábica en español:
% \usepackage[spanish,es-nodecimaldot]{babel}

% ---------- Geometría de página --------------------------------------
\usepackage[a4paper,margin=2.5cm]{geometry}

% ---------- Tablas ----------------------------------------------------
\usepackage{array}
\usepackage{makecell}
\usepackage[table]{xcolor}
\usepackage{ragged2e}
\usepackage{float}        % habilita el especificador [H] (tabla inquebrable)
\usepackage{graphicx}
\usepackage{caption}
\usepackage{fancyhdr}
\setlength{\headheight}{13.6pt}

% ---------- Enlaces ---------------------------------------------------
\usepackage{hyperref}
\hypersetup{
  colorlinks=true,
  linkcolor=black,
  citecolor=black,
  urlcolor=blue,
  pdftitle={Análisis DAR — Proxy-Caché y Almacenamiento de Archivos},
  pdfauthor={Santiago Quintero Uribe, Juan Manuel Perdomo Cárdenas, Andrés Felipe Zúñiga Zuluaga}
}

% ---------- Numeración de página estilo APA 7 ------------------------
% APA 7 para trabajos estudiantiles usa solo el número de página en el
% encabezado superior derecho. No se incluye running head salvo que se exija.
\pagestyle{fancy}
\fancyhf{}
\fancyhead[R]{\thepage}
\renewcommand{\headrulewidth}{0pt}

% ---------- Colores institucionales / de tabla ------------------------
\definecolor{gridhead}{HTML}{C5E0B3}   % verde claro de encabezado
\definecolor{gridrow}{HTML}{CCCCCC}    % gris de filas alternas
\definecolor{verdeuq}{HTML}{009B3A}    % verde Universidad del Quindío

% =====================================================================
%  LEYENDAS (captions) DE TABLAS
% =====================================================================
% position=top     -> el paquete caption calcula el espaciado asumiendo
%                     que el \caption va ENCIMA del contenido.
% skip=6pt         -> separación entre la leyenda y la tabla.
% justification    -> leyenda centrada si cabe en una línea, justificada
%                     si ocupa varias.
\captionsetup[table]{
  name=Tabla,
  labelfont=bf,
  labelsep=period,
  position=top,
  skip=6pt,
  font=small,
  justification=centering,
  singlelinecheck=true
}

% =====================================================================
%  APARIENCIA DE LAS TABLAS
% =====================================================================
\renewcommand{\arraystretch}{1.25}
\setlength{\tabcolsep}{4pt}

% Espacio antes/después de cada entorno table colocado con [H]
\setlength{\intextsep}{12pt plus 3pt minus 3pt}

% Tolerancia de página: evita que LaTeX estire en exceso el texto
% cuando una tabla se empuja a la página siguiente.
\raggedbottom
